\documentclass{article}

% Four-page, double-blind ML-for-Systems workshop extended abstract.

%\usepackage[dblblindworkshop]{neurips_2026}
\usepackage[dblblindworkshop]{neurips_2026}
\usepackage{comment}

\workshoptitle{ML for Systems Workshop}

\usepackage[utf8]{inputenc}

\usepackage[T1]{fontenc}

\usepackage{hyperref}

\usepackage{url}

\usepackage{booktabs}

\usepackage{graphicx}

\usepackage{amsmath}

\usepackage{amsfonts}

\usepackage{microtype}

\usepackage{xcolor}

% This is a planning document that compiles in the conference template. Replace

% blue text as evidence becomes available; do not submit with blueprint notes.

\definecolor{blueprint}{RGB}{25,82,145}

\newcommand{\todo}[1]{\textcolor{blueprint}{[TODO: #1]}}

\newcommand{\pilot}[1]{\textcolor{blueprint}{[PILOT: #1]}}

\newcommand{\bytes}[1]{\lvert #1\rvert_{\mathrm{UTF8}}}
\newcommand{\Dec}{\operatorname{D}}
\newcommand{\Tok}{\operatorname{T}}
\newcommand{\Ser}{\operatorname{SerializeCodec}}

\title{When Adversarial Prediction Still Compresses: Tokenization and Adversarial Robustness in LLM-Based Lossless Compression}

% Tokenization Buffers Prediction Failure in LLM-Based Lossless Compression

% Keep the submission anonymous. Add authors only in the camera-ready version.

\author{Anonymous Authors}

\begin{document}

\maketitle
Large language models (LLMs) have been known to be capable of acting as powerful lossless compression machines: next-token probabilities coupled with entropy coding can substantially compress data below its raw size. Yet existing evaluations largely single out predictive performance of the LLM or at best aggregate compression ratios, leaving an important systems question understudied: \emph{when does LLM compression survive (catastrophic/worst-case) prediction failure?} We show that the answer lies in treating LLM-based compression as a two-layer system of (1) tokenization as pre-compression and (2) entropy coding with LLM-derived probabilities. On text8, Qwen2.5-0.5B compresses to $17.6\%$ of raw size, outperforming token IDs ($41.3\%$) and Brotli ($29.6\%$). Strikingly, a $5.26\times$ increase in adversarial prediction cost still yields a $31.7\%$ arithmetic-coded stream because long tokens amortize prediction errors. Removing this effect with one-byte tokens causes expansion to $119.6\%$, while conventional codecs remain compressive. 
%Tokenizer fertility therefore defines a key robustness boundary, motivating selective predictive coding with a conventional fallback.



% -----------------------------------------------------------------------------

% BLUEPRINT AT A GLANCE (delete before submission)

% Target: at most four content pages, including figures and tables. Suggested

% budget: motivation 0.75; formulation and attacks 1.0; evaluation 0.75; results

% 0.9; systems implications, limitations, and conclusion 0.6. Do not assume

% references or supplementary material are excluded from the organizers four-page

% limit; confirm their counting rule before submission.

%

% One-sentence thesis:

% Tokenization is a variable-rate front end to predictive coding, so robust

% evaluation must optimize and report code length per decoded source byte, not

% merely next-token surprisal.

%

% Claim discipline:

% (C1) Formal: token surprisal and source-normalized compression ratio are

%      different objectives whenever decoded token lengths vary.

% (C2) Method: canonical, byte-aware greedy and beam attacks directly optimize

%      entropy-floor or realized arithmetic-code expansion.

% (C3) Empirical: use only after the full multi-model/seed matrix supports it.

%      The current repository establishes a Qwen2.5-0.5B case study, not broad

%      tokenizer or model universality.

% -----------------------------------------------------------------------------

\section{Who is actually doing the heavy lifting in LLM-based compression?}

% What is doing the compression?

% Proposed Story-Line: LLM-based lossless compressors predict tokens. Arithmetic coding converts token probabilities into bits. Performance is conventionally reported using NLL or perplexity. But users store bytes, not tokens. Tokens represent unequal byte lengths. Therefore, token-level worst cases may not be compression worst cases. Your experiment reveals exactly this mismatch.

``Language Modeling Is Compression'' introduced and formalized how next-symbol probabilities can drive a lossless arithmetic coder and mentions tokenization itself as lossless precompression~\citep{deletang2024language}. This observation implies that so-called LLM compression combines two distinct mechanisms: a tokenizer replaces variable-length source strings with token IDs, and a language model predicts those IDs so that an entropy coder can shorten likely ones. That the tokenization itself already provides some compressive benefit is often under-stated and has not been systematically, jointly with LLM-based entropy compression to detangle both compresssive effects studied. 

%LLMZip is another instance of this composed design~\citep{valmeekam2023llmzip}.

Prior work by Schmidt et al.~\citep{schmidt2026token} isolates the first mechanism and show that LLM token tables can themselves serve as global string dictionaries, enabling compression and operations over encoded data without model inference. Thus tokenization can be a useful compression primitive without language-model inference. This raises an attribution question for LLM codecs: how much compression comes from tokenization versus prediction?

%Schmidt et al.~\citep{schmidt2026token} isolate the first mechanism before introducing a ... . They use an LLM vocabulary as a global symbol table in which common byte strings map to fixed-width IDs and one-byte entries guarantee coverage. 

%Unlike a block-local code, this stable representation supports equality checks, hashing, joins, and aggregations before decoding. These properties make token tables particularly attractive for databases, where strings can remain encoded during common relational operations and be decoded only when needed. Thus tokenization can be a useful compression primitive without language-model inference. 

What remains unclear is how much of an LLM codec's behavior comes from this token representation and how much comes from prediction. Delétang et al. compare ASCII and BPE ~\citep{sennrich2016bpe} on natural enwik9, but the two effects have not been isolated under inputs chosen against realized end-to-end size. This matters because prediction loss is measured per token while compression is measured per source byte: a poor probability model may still sit inside a pipeline that compresses when its tokens represent many bytes, whereas low-fertility tokenization can expose expansion. We therefore study the tokenizer--predictor--coder pipeline as a two-layer system.

\paragraph{Contributions.}

We (i) factor end-to-end compression rate into tokenizer fertility and predictive coding under natural text, random input and worst-case inputs; (ii) we formalize token-, byte-, and realized-code objectives and evaluate against worst-case attacks per objective to demonstrate when the token layer absorbs predictive surporise and when it exposes expansion of the original file size and (iii) we propose a systems design for database compression that is LLM-native while still having the advantages of dictionary encodings (querabliity, ...). 

%We construct losslessly decodable adversarial inputs and separate the resulting file size into tokenizer, model, arithmetic-coder, and metadata costs.

\paragraph{Systems implication.}
The two-layer (tokenization + LLM-based entropy coding) view suggests a robust architecture for learned compression. Token IDs (global token table for all tables) can remain the stable, queryable representation used during processing, while predictive entropy coding exploiting LLM-next token probabilities is applied only as an optional storage, transport layer or to actually cold data. This separates the operational benefits of a global token representation from the inference cost and ... of the language model and introduces tokenization as LLM-native pre-compression. 
%to be determined: For near worst-case inputs, a codec can additionally store the smaller of the predictive representation and a conventional fallback, bounding expansion while retaining predictive gains on favorable inputs.


\section{Compression objectives and how to construct worst-case inputs for the predictor}
\label{sec:attacks}

%TODO: make introduction clearer. Right now, it feels like it comes out of nowwhere. Also, the writting is not good. 

Our central distinction is between prediction difficulty and end-to-end compression failure. We use one notation throughout. Let
$x=(x_1,\ldots,x_N)$ be the generated token sequence, let $h_t=x_{<t}$ be its token history, let
$\Tok$ and $\Dec$ denote the tokenizer and decoder, and let $p_\theta(v\mid h)$ be the language model probability of token $v$ after history $h$.
Define the decoded source size, in bytes, as
\begin{equation}
  B(x) = \bytes{\Dec(x)}.
  \label{eq:source-bytes}
\end{equation}
All rates below are normalized either by tokens or by these decoded source bytes.

For a candidate next token $v$, define its token surprisal
\begin{equation}
  \ell_\theta(v\mid h)
  = -\log_2 p_\theta(v\mid h)
  \qquad [\mathrm{bits/token}].
  \label{eq:token-surprisal}
\end{equation}
The token-level greedy attack selects the canonical continuation with maximum surprisal,
\begin{equation}
  v^\star_{\mathrm{tok}}
  = \arg\max_{v\in G(h)} \ell_\theta(v\mid h),
  \label{eq:minprob}
\end{equation}
where $G(h)$ is the set of admissible canonical next tokens. This is equivalent to selecting the minimum-probability continuation, but the score itself is written as surprisal so that its units are explicit.

To account for variable-length tokenization, define the incremental number of decoded source bytes contributed by $v$,
\begin{equation}
  \Delta B(v;h)
  =
  B(h\mathbin\Vert v)-B(h)
  =
  \bytes{\Dec(h\mathbin\Vert v)}-\bytes{\Dec(h)}.
  \label{eq:delta-bytes}
\end{equation}
The byte-aware greedy attack then selects
\begin{equation}
  v^\star_{\mathrm{byte}}
  =
  \arg\max_{v\in G(h)}
  \frac{\ell_\theta(v\mid h)}{\Delta B(v;h)}
  \qquad [\mathrm{bits/byte}].
  \label{eq:spb}
\end{equation}
Unlike Eq.~\ref{eq:minprob}, this objective directly targets ideal entropy-coding cost per newly decoded source byte.

For a complete sequence $x$, define the cumulative ideal model cost
\begin{equation}
  L_\theta(x)
  =
  \sum_{t=1}^{N}
  \ell_\theta(x_t\mid x_{<t})
  =
  -\sum_{t=1}^{N}\log_2 p_\theta(x_t\mid x_{<t})
  \qquad [\mathrm{bits}],
  \label{eq:nll}
\end{equation}
where any fixed initial context is implicit in $x_{<1}$. The ideal source-normalized rate and corresponding size ratio are
\begin{align}
  r_{\mathrm{ideal}}(x)
  &= \frac{L_\theta(x)}{B(x)}
  &&[\mathrm{bits/byte}], \label{eq:ideal-bpb}\\
  R_{\mathrm{ideal}}(x)
  &= \frac{L_\theta(x)}{8B(x)}
  =\frac{r_{\mathrm{ideal}}(x)}{8}
  &&[\text{dimensionless}]. \label{eq:ideal-ratio}
\end{align}
Thus $R_{\mathrm{ideal}}<1$ means that the ideal predictive code is smaller than raw 8-bit source bytes.

For the realized codec, let
\begin{equation}
  S_{\mathrm{codec}}(x)
  = \left|\Ser(x)\right|
  \qquad [\mathrm{bytes}]
  \label{eq:serialized-bytes}
\end{equation}
be the serialized output size. Its realized source-normalized rate and size ratio are
\begin{align}
  r_{\mathrm{real}}(x)
  &= \frac{8S_{\mathrm{codec}}(x)}{B(x)}
  &&[\mathrm{bits/byte}], \label{eq:real-bpb}\\
  R_{\mathrm{real}}(x)
  &= \frac{S_{\mathrm{codec}}(x)}{B(x)}
  =\frac{r_{\mathrm{real}}(x)}{8}
  &&[\text{dimensionless}]. \label{eq:actualrate}
\end{align}
A realized-size beam search approximately maximizes Eq.~\ref{eq:actualrate} by cloning coder state for candidate continuations. This objective captures arithmetic-coder, metadata, and framing effects that the ideal NLL does not.

\subsection{Why prediction loss is not compression loss}
\label{sec:objective-mismatch}

Define tokenizer fertility as
\begin{equation}
  f(x)
  = \frac{B(x)}{N}
  \qquad [\mathrm{bytes/token}].
  \label{eq:fertility}
\end{equation}
Combining Eqs.~\ref{eq:nll}, \ref{eq:ideal-bpb}, and \ref{eq:fertility} gives the exact factorization
\begin{equation}
  r_{\mathrm{ideal}}(x)
  =
  \underbrace{\frac{L_\theta(x)}{N}}_{\text{bits/token}}
  \underbrace{\frac{N}{B(x)}}_{\text{tokens/byte}}
  =
  \frac{L_\theta(x)/N}{f(x)}.
  \label{eq:factorization}
\end{equation}
Equivalently,
\begin{equation}
  R_{\mathrm{ideal}}(x)
  =
  \frac{L_\theta(x)/N}{8f(x)}.
  \label{eq:factorization-ratio}
\end{equation}
This makes the objective mismatch explicit: high prediction cost in bits/token can be offset by high tokenizer fertility in bytes/token.

Ignoring realized-coder and metadata overhead, predictive coding compresses relative to raw UTF-8 exactly when
\begin{equation}
  R_{\mathrm{ideal}}(x)<1
  \quad\Longleftrightarrow\quad
  \frac{L_\theta(x)}{N}<8f(x).
  \label{eq:compression-boundary}
\end{equation}
Hence a sequence with fertility $f$ bytes/token can tolerate an average prediction cost of up to $8f$ bits/token before the ideal predictive layer expands the source.

\subsection{Adversarial model and generation policies}

We assume white-box access to the compression pipeline: the adversary can inspect the tokenizer, model probabilities, and coding algorithm. It constructs a length-$N$ token sequence autoregressively. At history $h$, the admissible set $G(h)$ contains only continuations that preserve canonicality. Concretely, every generated prefix must satisfy
\begin{equation}
  \Tok\!\left(\Dec(x_{1:t})\right)=x_{1:t},
  \qquad t=1,\ldots,N,
  \label{eq:canonicality}
\end{equation}
so decoding and retokenization recover exactly the same token IDs and source bytes.

We compare three generation policies against natural text from Wikipedia (text8 in UTF-8) and random canonical sequences as baselines:
\emph{minimum probability} uses the token-level objective in Eq.~\ref{eq:minprob};
\emph{maximum surprisal/byte} uses Eq.~\ref{eq:spb}; and
\emph{realized-size beam} searches over multiple canonical continuations to approximately maximize the realized size ratio in Eq.~\ref{eq:actualrate}.

%\emph{Natural text} is the non-adversarial corpus baseline. %\emph{Random canonical} uniformly samples valid canonical extensions and controls for unusual token sequences without optimizing against the model.


%TODO: I do not yet know how to include this / where to best put it. 

%The \emph{one-byte constraint} is not a separate attack objective. It is an ablation of the candidate alphabet in which $G$ contains only canonical tokens contributing one decoded byte. This fixes fertility near one byte/token and removes the multi-byte-token buffering mechanism. Under an exact one-byte constraint, MinProb and SurprisalPerByte induce the same greedy ranking.

The \emph{one-byte constraint} restricts $G(h)$ to canonical tokens with $\Delta B(v;h)=1$. This removes variation in incremental token fertility, so Eqs.~\ref{eq:minprob} and \ref{eq:spb} induce the same greedy ranking.

%TODO: think about if this can be shorted or removed. If it stays, call the masking + rescoring an optimization of LLM-based compression. 
\paragraph{Global mask re-scoring as an optimization} 
When evaluating occurring-token dictionaries, we construct each adversarial sequence once under the full distribution and replay identical IDs and histories under the post-hoc occurring-token mask. This isolates probability renormalization; it is not a mask-adaptive attack.

\section{Experimental design}
\label{sec:design}

\paragraph{Setup.}
We evaluate Qwen2.5-0.5B on natural text from text8 and on canonical
adversarial sequences generated by the attacks in
Section~\ref{sec:attacks}. The natural-text baseline uses the first
100,000 text8 tokens; adversarial results are averaged over three
independently initialized 1,000-token sequences and reported as mean
$\pm$ sample standard deviation. All attacks use the same model,
context, and arithmetic-coding configuration. Full experimental
parameters and per-run statistics are reported in Appendix~\ref{app:experimental-details}.

\paragraph{Reported quantities.}
For every input or attack we report the same notation and units:
\begin{align}
  \text{Bits/token}
  &= \frac{L_\theta(x)}{N},
  \label{eq:reported-bpt}\\
  \text{Bytes/token}
  &= f(x)=\frac{B(x)}{N},
  \label{eq:reported-fertility}\\
  \text{Ideal bits/byte}
  &= r_{\mathrm{ideal}}(x)=\frac{L_\theta(x)}{B(x)},
  \label{eq:reported-ideal-bpb}\\
  \text{Actual bits/byte}
  &= r_{\mathrm{real}}(x)=\frac{8S_{\mathrm{codec}}(x)}{B(x)},
  \label{eq:reported-real-bpb}\\
  \text{Realized size ratio }R
  &= R_{\mathrm{real}}(x)=\frac{S_{\mathrm{codec}}(x)}{B(x)}.
  \label{eq:reported-ratio}
\end{align}
The first four quantities carry the units shown in their names; $R$ is dimensionless. In particular, $R<1$ denotes compression relative to raw UTF-8, whereas $R>1$ denotes expansion.

%Where a realized codec run has not yet been measured, we report ``--''.

\begin{table}[t]

\centering

\caption{
\textbf{Source-normalized predictive coding under increasingly adverse inputs.} Token-level surprisal / LLM prediction cost alone is not predictive of compression failure. Arrows indicate the direction favorable to compression ($\downarrow$ lower is better; $\uparrow$ higher is better). AC size $R_{\mathrm{AC}}$ reports the arithmetic-coder payload relative to raw UTF-8; $^\dagger$ Preliminary 32- and 128-token smoke tests only (three seed runs each).
$^\ddagger$ Preliminary 1,024-token smoke test (one seed run).
}

\label{tab:decisive}

\small
\setlength{\tabcolsep}{3.2pt}

\begin{tabular}{lrrrrr}

\toprule

\multicolumn{1}{l}{Input / condition} &
\multicolumn{1}{c}{Prediction cost $\downarrow$} &
\multicolumn{1}{c}{Tokenizer fertility $\uparrow$} &
\multicolumn{1}{c}{Ideal predictive rate $\downarrow$} &
\multicolumn{1}{c}{AC payload rate $\downarrow$} &
\multicolumn{1}{c}{AC size $R_{\mathrm{AC}}$ $\downarrow$} \\

\multicolumn{1}{l}{} &
\multicolumn{1}{c}{\footnotesize (bits/token)} &
\multicolumn{1}{c}{\footnotesize (bytes/token)} &
\multicolumn{1}{c}{\footnotesize (bits/byte)} &
\multicolumn{1}{c}{\footnotesize (bits/byte)} &
\multicolumn{1}{c}{\footnotesize (\% of raw)} \\

\midrule

\multicolumn{6}{l}{\textit{Baselines}} \\

Natural text (text8, UTF-8)
    & 5.11
    & 5.59
    & 0.92
    & 1.13
    & $14{,}1\%$ \\

Random canonical (UTF-8)
    & 8.74
    & 1.30
    & 6.71
    & 7.59
    & $94{,}8\%$ \\

\midrule

\multicolumn{6}{l}{\textit{Adversaries}} \\

MinProb adversary
    & $26.86 \pm 0.05$
    & $7.09 \pm 0.04$
    & $3.79 \pm 0.02$
    & $2.53 \pm 0.01$
    & $31{,}7 \pm 0{,}2\%$ \\

Max. surprisal/byte ($N=32$)$^\dagger$
    & $19.26 \pm 0.39$
    & $1.45 \pm 0.02$
    & $12.89 \pm 0.40$
    & $10.74 \pm 0.31$
    & $134{,}2 \pm 3{,}9\%$ \\

Max. surprisal/byte ($N=128$)$^\dagger$
    & $17.82 \pm 0.13$
    & $1.48 \pm 0.01$
    & $11.96 \pm 0.07$
    & $10.43 \pm 0.06$
    & $130{,}4 \pm 0{,}7\%$ \\

Max. surprisal/byte ($N=1024$)$^\ddagger$
    & $17.49 \pm 0.00$
    & $1.47 \pm 0.00$
    & $11.87 \pm 0.00$
    & $10.50 \pm 0.00$
    & $131{,}3 \pm 0{,}0\%$ \\

Realized-size beam
    & --- & --- & --- & --- & --- \\

\midrule

\multicolumn{6}{l}{\textit{Tokenizer ablation}} \\

One-byte UTF-8 ablation
    & $8.33 \pm 0.00$
    & $1.00 \pm 0.00$
    & $8.33 \pm 0.00$
    & $9.57 \pm 0.00$
    & $119{,}6 \pm 0{,}0\%$ \\

\bottomrule

\end{tabular}

\end{table}

\section{Results}
\label{sec:results}

\begin{figure}[t]
    \centering
    \includegraphics[width=\linewidth]{plots/end_to_end_compressor_robustness_preview.png}
    \caption{\textbf{End-to-end compression across natural and adversarial inputs by codec.}\\ Realized size ratio $R$ across baseline and adversarial inputs for the codecs Qwen + AC, tokenization only using Qwens tokenizer, FSST, and brotli. Qwen tokenization stores fixed-width token IDs. Bold values mark the best result per input.}
    \label{fig:layered_compression_robustness}
\end{figure}


\paragraph{Prediction failure does not imply compression failure.}
The existing minimum-probability full-vocabulary attack reaches $28.409\pm0.077$ ideal bits per predicted token, $5.734\times$ the matched text8 value of 4.954 bits/token. Because a fixed ID for the 151,643-token vocabulary requires 18 bits, prediction makes this adversarial token stream substantially more expensive than fixed-width IDs. Yet the sequences decode to an average 6.914 source bytes/token. Dividing model cost by this fertility gives approximately 4.108 ideal bits/source byte, or 51.346\% of raw UTF-8. The model therefore fails badly at the token level while the composed tokenizer--predictor system still has $R_{\mathrm{ideal}}<1$.

\paragraph{Byte-aware attacks test the relevant compression objective.}
Minimum probability maximizes token surprisal without accounting for $\Delta B$. Maximum surprisal/byte optimizes the local source-normalized ideal rate, while realized-size beam targets the serialized size ratio $R_{\mathrm{real}}$ directly. Table~\ref{tab:decisive} is therefore the primary comparison: it determines whether progressively more end-to-end objectives increase ideal or realized bits/source byte. The corresponding full-vocabulary SurprisalPerByte and realized- size beam measurements are required before making a claim that MinProb is the strongest compression attack.

\paragraph{Removing the tokenization gain exposes predictive expansion.}

The one-byte-ASCII ablation fixes fertility at one byte/token and removes the tokenizer's ability to compensate for poor prediction using long tokens. In this setting, every accepted token contributes exactly one decoded byte, fixing fertility at one byte/token. Consequently, MinProb and SurprisalPerByte become identical and the break-even model cost is exactly eight bits/token.

The measured model floor is 135.48\% of raw bytes and the arithmetic payload is 143.36\%; an NLL beam reaches 136.70\% ideal and 144.14\% arithmetic payload. Thus the predictive layer cannot beat the original 8-bit bytes once the tokenizer is prevented from absorbing model surprise through multi-byte tokens. Thus, once first-stage token compression is removed, the same
kind of prediction failure that remained hidden by multi-byte Qwen tokens becomes direct end-to-end expansion. Mind that the narrow 32-token beam is an approximation of worst-case input.

\begin{comment}
\paragraph{Dictionary masking changes prediction cost, not the attack sequence.}
Post-hoc occurring-token dictionaries contain 592, 606, and 587 symbols for the three existing adversarial sequences. Replaying identical IDs and histories under these masks lowers adversarial cost from $28.409\pm0.077$ to $18.402\pm0.833$ bits/token, a 35.2\% reduction, and lowers the byte-normalized ideal rate from 51.346\% to 33.284\%. It remains $3.806\times$ the matched masked- text8 cost of 4.835 bits/token. Because generation is unchanged, this experiment isolates renormalization rather than adaptive robustness.
\end{comment}

\paragraph{Realized overhead must be reported separately from the ideal floor.}
For text8, model-induced floors are 11.631\% of raw UTF-8 with the full vocabulary and 11.352\% with occurring tokens. Finite coding and metadata raise realized totals to 14.405\% and 13.434\%, corresponding to $6.942\times$ and $7.444\times$ compression. The occurring-token bitmap costs 10,818 bytes (2.03\% of raw size), less than its coding gain at 100,000 tokens. These gaps motivate keeping ideal and actual columns separate in Table~\ref{tab:decisive}.

\section{Systems implications and limitations}

\label{sec:limitations}

The systems lesson is to preserve token IDs as the global, queryable in-memory representation and treat predictive arithmetic coding as an optional block-level storage or transport layer. 

Source-normalized accounting identifies whether fragility comes from the tokenizer, predictor, quantizer, dictionary, or container.

For high-cost inputs blocks (meaning ones with data-model mismatch), a guard can store the smaller of predictive and conventional encodings plus a selector bit, bounding size by the fallback plus one bit; its decoding granularity, latency, and throughput remain to be measured.


%TODO: rephrase this limitations paragraph
%Our evidence is limited to one model, synthetic white-box inputs, three starts, finite context, approximate search, and one coder/container. It does not establish global worst cases or generalize to multilingual text or code. The attack could enable storage or compute amplification; practical mitigations include size caps, early-abort expansion detection, bounded inference, and fallback codecs.

\section{Conclusion}

Robust learned compression is an end-to-end systems problem. Canonical byte-aware attacks and explicit cost attribution show that the tested adversarial sequences are 3.81--5.73 times harder than matched text8, while guarded fallback offers a principled path from diagnosis to robust operation.

\begin{ack}

\todo{Camera-ready only: funding and competing-interest disclosures. The supplied style hides this environment in anonymous submission mode.}

\end{ack}

\bibliographystyle{plainnat} % use the style required by the workshop/template
\bibliography{references}


\newpage
\section{Appendix}


\paragraph{Why canonicality is required.}
Not every sequence of valid token IDs can arise from tokenizing its own decoded
bytes. For example, suppose the vocabulary contains tokens for
\texttt{a}, \texttt{b}, and \texttt{ab}. The token sequence
$[\texttt{a},\texttt{b}]$ decodes to \texttt{ab}, but the tokenizer may encode
\texttt{ab} as the single token $[\texttt{ab}]$. Hence, decoding and
retokenizing changes the token sequence. Since our attacks select token IDs
directly from the model distribution, allowing such sequences would optimize
over token histories that cannot occur for the corresponding source input.
We therefore require every generated partial sequence to be canonical, i.e.,
$T(\operatorname{decode}(x_{1:t}))=x_{1:t}$.



\subsection{Experimental details}
\label{app:experimental-details}

\paragraph{Model and coder.}
We use Qwen2.5-0.5B with 151,643 non-special candidate tokens, a context length of 1,000 tokens, 100 retained context tokens between coding blocks, and KV caching. The arithmetic coder uses a 32-bit state and a frequency total of 262,144. The finalized QACS version-1 stream records the seed, model and coder configuration, token count, raw byte count, arithmetic payload, any required token dictionary, and framing. Model weights and the shared tokenizer table are treated as shared decoder state and excluded from both Qwen representations.

\paragraph{Inputs and aggregation.}
Every production sequence in the primary attack evaluation contains 10,000 tokens. Table~\ref{tab:decisive} reports natural text, random canonical input, the full-vocabulary minimum-probability attack, and the one-byte UTF-8 ablation under this common token budget; the maximum-surprisal-per-byte row is populated only by a marked 32-token smoke test, and the missing realized-size measurement is left blank. Attack statistics are means and sample standard deviations over three independent canonical seed runs.

Figure~\ref{fig:layered_compression_robustness} and Appendix Table~\ref{tab:codec-comparison} are a separate codec-level diagnostic. For that comparison only, we serialize representative prefixes with 9,996--10,000 decoded source bytes so that complete-stream sizes are directly comparable across codecs. Thus Table~\ref{tab:decisive} reports arithmetic payload ratio $R_{\mathrm{AC}}$ for the token-budget runs, whereas the figure reports complete-stream $R_{\mathrm{real}}$ for byte-matched prefixes.

All generated prefixes satisfy
\[
\Tok(\Dec(x_{1:t}))=x_{1:t}, \qquad t=1,\ldots,N,
\]
and every final stream is independently decoded back to identical token IDs and UTF-8 bytes. FSST uses the official implementation at commit \texttt{e638d4cf8c26129d73c242a4127b42b975de5b63}; Brotli uses quality~11; token IDs use a fixed 18-bit representation with a stream header.

\paragraph{Cost accounting.}
In the auxiliary codec comparison, we report both the Qwen+AC arithmetic payload and the complete serialized stream. Their difference is 331--334 bytes across the five byte-matched inputs. The finite-frequency table reserves a nonzero count for every vocabulary entry, so extremely small model probabilities are clipped. Consequently, realized arithmetic length can be lower than unquantized NLL on the minimum-probability attack; this quantization effect is reported rather than interpreted as predictive accuracy.

\paragraph{Layered compression robustness and cost attribution.} \textbf{(a)} End-to-end size relative to raw UTF-8, separating fixed-width token IDs, the ideal model surprisal cost, the realized arithmetic-coding (AC) payload plus dictionary metadata, and the final serialized file (dot). The serialized size can be decomposed as ($S_{\mathrm{total}} = S_{\mathrm{payload}} + S_{\mathrm{dict}} + S_{\mathrm{frame}}$), where $S_{\mathrm{payload}}$ is the finalized arithmetic payload, $S_{\mathrm{dict}}$ is dictionary metadata, and $S_{\mathrm{frame}}$ is the remaining container and framing overhead.

\begin{table}[t]
\centering
\caption{\textbf{Realized size ratio by input and codec.} Values are encoded bytes divided by raw UTF-8 bytes and reported as percentages; lower is better. AC payload excludes the Qwen container and is shown only for cost attribution. All entries in the four codec columns are complete, round-trip-verified streams, and bold marks the smallest complete stream per row.}
\label{tab:codec-comparison}
\small
\setlength{\tabcolsep}{3.4pt}
\begin{tabular}{lrrrrr}
\toprule
Input & AC payload & Qwen+AC & Token IDs & FSST & Brotli \\
\midrule
text8                    & $14.3\%$ & \textbf{$17.6\%$} & $41.3\%$ & $56.3\%$ & $29.6\%$ \\
Random printable         & $94.8\%$ & $97.4\%$ & $172.1\%$ & $98.2\%$ & \textbf{$83.0\%$} \\
MinProb ($\vtok$) adversary         & $32.9\%$ & $36.2\%$ & \textbf{$33.2\%$} & $79.5\%$ & $46.7\%$ \\
MaxSurprisal/Byte ($\vbyte$; $N=1{,}024$) & $122.4\%$ & $144.3\%$ & $154.4\%$ & $94.6\%$ & \textbf{$83.2\%$} \\
One-byte UTF-8 ablation & $119.5\%$ & $122.8\%$ & $225.2\%$ & $83.6\%$ & \textbf{$76.1\%$} \\
\bottomrule
\end{tabular}
\end{table}


The printable-only row is retained solely as an alphabet-sensitivity check: restricting the candidate set to the 95 printable ASCII bytes weakens, but does not remove, predictive expansion.


\paragraph{Tokenizer fertility stats}

\begin{table}[t]
  \centering
  \small
  \caption{Tokenizer fertility on the held-out 5\,MB text8 test region
  (855{,}233 whitespace-delimited words). The text8 BPE is trained only on the
  first 90\,MB with a 32k vocabulary. Fertility is tokens per word (lower is
  better); bytes per token is the compression-oriented inverse (higher is
  better). All tokenizers reconstruct the input exactly. The domain-trained BPE
  produces 4.4\% fewer tokens than Qwen.}
  \label{tab:tokenizer-fertility}
  \begin{tabular}{@{}lrrrrr@{}}
    \toprule
    Tokenizer & Vocab. & Used & Tokens & Tok./word $\downarrow$ & Bytes/tok. $\uparrow$ \\
    \midrule
    ASCII byte & 128 & 27 & 5{,}000{,}000 & 5.8464 & 1.0000 \\
    text8 BPE (32k) & 32{,}000 & 26{,}617 & 933{,}499 & \textbf{1.0915} & \textbf{5.3562} \\
    Qwen2.5 & 151{,}665 & 25{,}798 & 976{,}828 & 1.1422 & 5.1186 \\
    \bottomrule
  \end{tabular}
\end{table}




\paragraph{Reported quantities.}
The primary and auxiliary evaluations share the same token- and source-normalized quantities but differ in stream accounting:

\begin{align}
  \text{Bits/token} &= \frac{L_\theta(x)}{N},\\
  \text{Bytes/token} &= \frac{\bytes{\operatorname{decode}(x)}}{N},\\
  \text{Ideal bits/byte} &= \frac{L_\theta(x)}
    {\bytes{\operatorname{decode}(x)}},\\
  \text{AC bits/byte} &= \frac{S_{\mathrm{AC}}(x)}
    {\bytes{\operatorname{decode}(x)}},\\
  R_{\mathrm{AC}} &= 100\,\frac{S_{\mathrm{AC}}(x)}
    {8\,\bytes{\operatorname{decode}(x)}},\\
  R_{\mathrm{real}} &= 100\,\frac{\left|\operatorname{SerializeCodec}(x)\right|}
    {\bytes{\operatorname{decode}(x)}} \quad [\%].
\end{align}

Both ratios are relative to raw UTF-8; $R<100\%$ denotes compression and $R>100\%$ denotes expansion.


% Supplementary planning notes (do not count toward the four-page narrative):

% - Prove the token/byte objective mismatch and give attack pseudocode.

% - Record configurations, revisions, commands, seeds, hardware, compute, and

%   licenses; report full per-run tables and search-budget ablations.

% - Document exact seed-token, quantization, bitmap, payload, and header costs.

% - Repository map: src/adversarial.py; src/compression\_attacks.py;

%   scripts/run\_compression\_attacks.py; evaluation/notebooks/neurips/

%   adversarial\_worst\_case.ipynb; and artifacts/runs/adversarial/

%   qwen\_05b\_ascii\_bytes\_n1000.

% - Confirm whether references and the NeurIPS checklist are outside the

%   organizers four-page limit. Keep the checklist unless organizers explicitly

%   say their workshop does not require it.


%TODOS:

% put idea in there that tokenization is base-compression and then llm based compression can be used on top of that. This ensures that we can still query some of the content (please read tokenization compression paper again).

% include classical baselines such as brotli, llmzip

% include tokenization only as baseline and report for different tokenizers (Qwen’s large vocabulary;, a smaller BPE tokenizer; ideally a byte-level or byte-fallback baseline)

% mention FSST and include it as baseline 

% cite the econimics of LLM based compression from Andreas Kipf 



\end{document}
